%==============================================================
\section{Biên dịch mạch lượng tử}
\label{sec:compilation}
%==============================================================

Các chương trước coi mạch lượng tử là một đối tượng toán học: viết \code{qc.ccx(0,1,2)} thì có cổng
Toffoli, viết \code{qc.cx(0,4)} thì có CX giữa qubit 0 và 4. Phần cứng thật không cho phép điều đó.
Chương này phân tích quá trình \textbf{compilation} -- dịch mạch lý tưởng thành mạch mà một thiết bị
cụ thể chạy được, với sai số nhỏ nhất có thể -- và nhìn nó như một \textbf{bài toán tối ưu}.

\subsection{Phát biểu bài toán}

\begin{definition}[Quantum circuit compilation]
Cho mạch $C$ trên $n$ qubit logic, và một thiết bị mô tả bởi bộ ba
$\mathcal{T} = (\mathcal{G}, G_{\text{cm}}, \varepsilon)$, trong đó $\mathcal{G}$ là tập
\textbf{basis gates}, $G_{\text{cm}} = (V, E)$ là \textbf{coupling map} (đỉnh là qubit vật lý, cạnh
là cặp qubit cho phép cổng hai qubit) và $\varepsilon$ gán sai số cho từng cổng trên từng qubit/cạnh.
Tìm mạch $C'$ sao cho: (i) $C'$ chỉ dùng cổng trong $\mathcal{G}$; (ii) mọi cổng hai qubit của $C'$
nằm trên một cạnh của $G_{\text{cm}}$; (iii) $C'$ tương đương với $C$ (cùng unitary, sai khác một
hoán vị qubit đã biết và global phase); và (iv) $C'$ tối ưu theo một hàm cost (Mục~\ref{sec:cost}).
\end{definition}

Trong Qiskit, bộ ba $\mathcal{T}$ được đóng gói thành đối tượng \code{Target}. Điều kiện (iii) là chỗ
dễ sai nhất khi tự viết compiler, và Mục~\ref{sec:verify} sẽ chỉ ra vì sao.

\subsection{Pipeline của transpiler}

Qiskit không biên dịch bằng một hàm lớn duy nhất mà chạy một chuỗi \textbf{pass}, mỗi pass làm một
việc nhỏ. Các pass được nhóm thành sáu \textbf{stage} cố định (Hình~\ref{fig:pipeline},
Bảng~\ref{tab:stages}); hàm \code{generate\_preset\_pass\_manager} dựng sẵn chuỗi này theo tham số
\code{optimization\_level} và trả về một \code{StagedPassManager}.

\begin{figure}[h]
  \centering
  \begin{tikzpicture}[node distance=0.35cm, every node/.style={font=\small}]
    \node[box=black, minimum width=1.7cm] (in) {\code{QuantumCircuit}};
    \node[box=AccentBlue, right=of in, minimum width=1.2cm] (init) {\code{init}};
    \node[box=OkGreen, right=of init, minimum width=1.2cm] (lay) {\code{layout}};
    \node[box=OkGreen, right=of lay, minimum width=1.2cm] (rou) {\code{routing}};
    \node[box=OrangeYellow, right=of rou, minimum width=1.2cm] (tra) {\code{translation}};
    \node[box=Purple, right=of tra, minimum width=1.2cm] (opt) {\code{optimization}};
    \node[box=AccentBlue, below=0.9cm of opt, minimum width=1.2cm] (sch) {\code{scheduling}};
    \node[box=ProbRed, left=of sch, minimum width=1.7cm] (out) {mạch chạy được};
    \draw[arr] (in) -- (init);
    \draw[arr] (init) -- (lay);
    \draw[arr] (lay) -- (rou);
    \draw[arr] (rou) -- (tra);
    \draw[arr] (tra) -- (opt);
    \draw[arr] (opt) -- (sch);
    \draw[arr] (sch) -- (out);
    \draw[arr, dashed] (opt.north east) to[out=70, in=110, looseness=1.6] node[above, font=\footnotesize] {lặp tới khi hội tụ} (opt.north west);
    \node[draw=black!25, dashed, rounded corners, fit=(lay)(rou), inner sep=4pt,
          label={[font=\footnotesize, text=OkGreen]below:connectivity}] {};
  \end{tikzpicture}
  \caption{Sáu stage của pass manager trong Qiskit. Layout và routing giải quyết ràng buộc
  connectivity; translation giải quyết ràng buộc basis gates; optimization lặp đến khi không giảm
  được nữa.}
  \label{fig:pipeline}
\end{figure}

\begin{table}[h]
  \centering
  \caption{Sáu stage của \code{StagedPassManager} và các pass tiêu biểu.}
  \label{tab:stages}
  \small
  \begin{tabularx}{\linewidth}{@{}lXl@{}}
    \toprule
    \textbf{Stage} & \textbf{Nhiệm vụ} & \textbf{Pass tiêu biểu} \\
    \midrule
    \code{init} & chuẩn hoá đầu vào, phân rã cổng $\ge 3$ qubit & \code{HighLevelSynthesis},\\ & & \code{UnitarySynthesis} \\
    \code{layout} & chọn qubit vật lý cho từng qubit logic & \code{SabreLayout}, \code{VF2Layout} \\
    \code{routing} & chèn SWAP để mọi cổng hai qubit hợp lệ & \code{SabreSwap}, \code{BasicSwap} \\
    \code{translation} & dịch mọi cổng về basis gates & \code{BasisTranslator} \\
    \code{optimization} & huỷ, gộp cổng; lặp tới khi hội tụ & \code{Optimize1qGatesDecomposition},\\ & & \code{CommutativeCancellation} \\
    \code{scheduling} & gán thời điểm thực thi, chèn delay & \code{ALAPScheduleAnalysis}, \code{PadDelay} \\
    \bottomrule
  \end{tabularx}
\end{table}

Theo dõi từng pass bằng tham số \code{callback} của \code{pm.run} cho thấy, với mạch ví dụ ở
Hình~\ref{fig:motivation}, \textbf{47 pass} được chạy ở \code{optimization\_level=2}.
\code{BasisTranslator} là pass làm mạch \emph{nở ra} nhiều nhất (dịch \code{ccx}, \code{ry},
\code{cz} ra basis), các pass optimization sau đó \emph{co} mạch lại, còn phần lớn pass không đổi
mạch -- chúng là pass phân tích. Có hai loại pass:
\begin{itemize}[noitemsep]
  \item \code{AnalysisPass}: đọc mạch, ghi kết quả vào \code{property\_set} -- một từ điển dùng
    chung cho cả pipeline và là kênh duy nhất để các pass trao đổi thông tin; \textbf{không} sửa mạch.
  \item \code{TransformationPass}: nhận mạch, trả về mạch mới.
\end{itemize}

\subsection{\code{DAGCircuit} và front layer}

\code{QuantumCircuit} là một \emph{danh sách} cổng theo thứ tự viết. Thứ tự đó giả tạo: hai cổng trên
hai qubit khác nhau viết trước--sau nhưng thực ra độc lập. Vì vậy pass làm việc trên biểu diễn trung
gian \code{DAGCircuit} -- một \textbf{đồ thị có hướng không chu trình} (directed acyclic graph) mà
đỉnh là cổng và cạnh là qubit/clbit nối hai cổng liên tiếp. Cạnh biểu diễn quan hệ phụ thuộc thật:
cổng $B$ nằm sau cổng $A$ trên DAG nghĩa là $B$ \emph{phải} chờ $A$.

\begin{definition}[Front layer]
\textbf{Front layer} $F$ của một DAG là tập các cổng không còn cổng nào chưa thực thi đứng trước nó
-- tức tập cổng \emph{có thể thực thi ngay}.
\end{definition}

Front layer là khái niệm trung tâm của routing (Mục~\ref{sec:routing}) và chính là thành phần chính
của \emph{state} khi phát biểu routing thành bài toán MDP (Mục~\ref{sec:mdp}).

\subsection{Hàm cost: thế nào là một mạch ``tốt''?}
\label{sec:cost}

Có bốn thước đo hay dùng, và chúng \textbf{không cùng chiều}:
\begin{itemize}[noitemsep]
  \item \textbf{depth} -- số lớp cổng trên đường dài nhất; tỉ lệ với thời gian chạy, cần so với thời
    gian decoherence $T_2$.
  \item \textbf{size} -- tổng số cổng.
  \item \textbf{số cổng hai qubit} -- chi phối sai số, nên có trọng số lớn nhất.
  \item \textbf{số SWAP} -- giá phải trả cho connectivity; mỗi SWAP bằng 3 CX.
\end{itemize}
Một ước lượng thô nhưng hữu dụng: nếu mạch có $m$ cổng hai qubit, mỗi cổng lỗi với xác suất
$\epsilon$, thì xác suất chạy trọn mạch không lỗi là
\begin{equation}
  P_{\text{thành công}} \approx (1-\epsilon)^m \approx e^{-\epsilon m}, \qquad
  \text{hoặc chính xác hơn}\quad P_{\text{thành công}} \approx \prod_{g} (1 - \epsilon_g),
  \label{eq:psuccess}
\end{equation}
trong đó tích lấy trên mọi cổng $g$ với sai số $\epsilon_g$ riêng của qubit/cạnh mà nó chạy trên
đó (đọc từ \code{Target}). Bảng~\ref{tab:optlevel} cho thấy tác động của compiler: cùng mạch, cùng
thiết bị giả lập (\code{GenericBackendV2}, sai số \cx{} trung bình $3.1\times10^{-3}$), chỉ đổi
\code{optimization\_level} đã nâng xác suất thành công từ $0.89$ lên $0.94$.

\begin{table}[h]
  \centering
  \caption{Ảnh hưởng của \code{optimization\_level} lên mạch ví dụ (5 qubit, coupling map
  \code{line(5)}). $P$ tính theo~\eqref{eq:psuccess} chỉ với các cổng \cx{}.}
  \label{tab:optlevel}
  \small
  \begin{tabular}{@{}cccccc@{}}
    \toprule
    \code{optimization\_level} & depth & size & số \cx{} & cổng 1 qubit & $P_{\text{thành công}}$ \\
    \midrule
    0 & 57 & 82 & 37 & 40 & 0.892 \\
    1 & 46 & 67 & 31 & 31 & 0.908 \\
    2 & 40 & 65 & \textbf{19} & 41 & \textbf{0.943} \\
    3 & 40 & 65 & 19 & 41 & 0.943 \\
    \bottomrule
  \end{tabular}
\end{table}

Ba bài toán con quan trọng nhất của compiler -- \emph{layout}, \emph{routing} và \emph{synthesis} --
được trình bày lần lượt dưới đây.

%--------------------------------------------------------------
\subsection{Layout: gán qubit logic cho qubit vật lý}
\label{sec:layout}

Mạch dùng các \textbf{qubit logic} $q_0, q_1, \dots$ -- chỉ là nhãn. Thiết bị có các \textbf{qubit
vật lý} $Q_0, Q_1, \dots$ nối với nhau theo coupling map. \textbf{Layout} (còn gọi là
\emph{initial mapping} hay \emph{qubit allocation}) là một đơn ánh
$\pi: \{q_i\} \to \{Q_j\}$.

\paragraph{Coupling map và distance matrix.} Bậc trung bình của coupling map đo độ ``chật'' của
thiết bị: \code{line(5)} có bậc 1.6, \code{ring(5)} 2.0, \code{grid(2,4)} 2.5, \code{full(5)} 4.0;
kiến trúc heavy-hex của IBM chỉ khoảng 2.1. Dữ liệu quan trọng nhất của coupling map là
\textbf{distance matrix} $D$, với $D[i][j]$ là độ dài đường đi ngắn nhất giữa $Q_i$ và $Q_j$; cổng
hai qubit giữa $Q_i$ và $Q_j$ cần ít nhất $D[i][j]-1$ SWAP. Mọi heuristic routing đều tính cost từ $D$.

\paragraph{Layout là bài toán subgraph isomorphism.} Mỗi mạch có một \textbf{interaction graph}:
đỉnh là qubit logic, cạnh là cặp qubit có ít nhất một cổng hai qubit. Nếu interaction graph
\emph{nhúng được} vào coupling map, tồn tại layout cần \textbf{0 SWAP}. Tìm phép nhúng đó là bài toán
\emph{subgraph isomorphism} -- NP-đầy đủ; khi không nhúng được, tối thiểu hoá số SWAP cũng là
NP-hard \cite{siraichi2018allocation}.

\begin{table}[h]
  \centering
  \caption{Bốn layout pass của Qiskit, và kết quả trên cùng một mạch 4 qubit có 3 cổng \cx{}
  trên coupling map \code{line(4)}.}
  \label{tab:layout}
  \small
  \begin{tabular}{@{}llcc@{}}
    \toprule
    \textbf{Pass} & \textbf{Cách làm} & \textbf{SWAP} & \textbf{depth} \\
    \midrule
    \code{TrivialLayout} & $q_i \to Q_i$, làm baseline & 3 & 5 \\
    \code{DenseLayout} & chọn vùng con dày nhất của thiết bị, không nhìn mạch & 3 & 5 \\
    \code{VF2Layout} & subgraph isomorphism chính xác (VF2++ \cite{juttner2018vf2pp}) & \textbf{0} & \textbf{3} \\
    \code{SabreLayout} & heuristic SABRE, làm layout \emph{và} routing nhiều lần thử & -- & -- \\
    \bottomrule
  \end{tabular}
\end{table}

Bảng~\ref{tab:layout} cho thấy giá trị của layout: cùng một mạch, chỉ khác nhau ở cách đánh số qubit,
mà một bên cần 3 SWAP (tức 9 cổng \cx{} thêm vào một mạch vốn chỉ có 3 cổng \cx{}), bên kia cần 0.

\paragraph{Noise-aware layout.} Các cạnh của coupling map không như nhau: trên thiết bị giả lập, cạnh
tệ nhất có sai số gấp 4.1 lần cạnh tốt nhất. Pass \code{VF2PostLayout} tìm một phép nhúng khác cùng
hình dạng nhưng trên các cạnh tốt hơn. Trong thí nghiệm, nó giảm tổng sai số của một mạch 3 cổng
\cx{} khoảng \textbf{56\%} (xác suất không lỗi từ $0.9884$ lên $0.9949$) mà \emph{không bớt cổng nào}.

%--------------------------------------------------------------
\subsection{Routing: chèn SWAP}
\label{sec:routing}

Với mạch thực tế, interaction graph luôn dày hơn coupling map nên không layout nào đủ tốt; compiler
buộc phải chèn SWAP để đưa các cặp qubit cần tương tác lại cạnh nhau. Mọi router -- từ
\code{BasicSwap}, tới SABRE, tới một RL policy -- đều là cùng một vòng lặp (Thuật toán~\ref{alg:routing}),
chỉ khác nhau ở \textbf{một dòng}: chọn SWAP nào.

\begin{algorithm}[h]
  \caption{Vòng lặp routing tổng quát}
  \label{alg:routing}
  \begin{algorithmic}[1]
    \Require DAG của mạch, coupling map $G_{\text{cm}}=(V,E)$, layout ban đầu $\pi$
    \Ensure mạch hợp lệ với $G_{\text{cm}}$ và layout cuối $\pi$
    \State $F \gets$ front layer của DAG
    \While{DAG còn cổng}
      \State $\mathcal{E} \gets \{g \in F : g$ một qubit, hoặc $(\pi(g_1),\pi(g_2)) \in E\}$
      \If{$\mathcal{E} \neq \emptyset$}
        \State thực thi mọi cổng trong $\mathcal{E}$, xoá khỏi DAG, cập nhật $F$
      \Else
        \State chọn cạnh $(a,b) \in E$ \Comment{\textcolor{ProbRed}{quyết định duy nhất}}
        \State chèn $\mathrm{SWAP}(a,b)$, cập nhật $\pi$
      \EndIf
    \EndWhile
  \end{algorithmic}
\end{algorithm}

Tối thiểu hoá số SWAP là NP-hard, nên mọi router thực tế dùng heuristic. Thuật toán \textbf{SABRE}
\cite{li2019sabre} -- router mặc định của Qiskit (\code{SabreSwap}) -- chọn SWAP làm giảm nhiều nhất
một hàm cost $H$ tính trên distance matrix. Gọi $F$ là front layer và $E_{\text{ext}}$ là
\emph{extended set} (vài lớp cổng hai qubit kế tiếp), SABRE có ba heuristic:
\begin{align}
  H_{\text{basic}} &= \frac{1}{|F|}\sum_{g\in F} D\big[\pi(g_1)\big]\big[\pi(g_2)\big], \label{eq:hbasic}\\
  H_{\text{lookahead}} &= H_{\text{basic}} + W\cdot\frac{1}{|E_{\text{ext}}|}\sum_{g\in E_{\text{ext}}} D\big[\pi(g_1)\big]\big[\pi(g_2)\big], \label{eq:hlook}\\
  H_{\text{decay}} &= \max\big(\mathrm{decay}[a], \mathrm{decay}[b]\big)\cdot H_{\text{lookahead}}, \label{eq:hdecay}
\end{align}
trong đó $(a,b)$ là cạnh ứng viên và $\mathrm{decay}[\cdot]$ tăng mỗi khi một qubit vừa bị SWAP, để
tránh dồn mọi SWAP vào cùng một qubit làm depth tăng vọt. \code{decay} là mặc định, và đây là hàm
mục tiêu mà một RL policy cần vượt qua -- lưu ý nó đã cân bằng giữa số SWAP và depth chứ không chỉ
đếm SWAP. Ngoài ra, \code{SabreLayout} chạy routing \emph{xuôi rồi ngược} nhiều lần để tinh chỉnh
layout ban đầu, nên layout và routing trong SABRE không tách rời.

\begin{table}[h]
  \centering
  \caption{So sánh các router trên một mạch 5 qubit (29 \cx{}, depth 71), coupling map
  \code{line(5)}, layout tầm thường. \code{mini\_router} là router tự cài đặt theo đúng công
  thức~\eqref{eq:hbasic}.}
  \label{tab:routers}
  \small
  \begin{tabular}{@{}lcc@{}}
    \toprule
    \textbf{Router} & \textbf{SWAP} & \textbf{depth} \\
    \midrule
    \code{BasicSwap} (đường ngắn nhất, không nhìn trước) & 14 & 123 \\
    \code{LookaheadSwap} & 7 & 89 \\
    \code{SabreSwap} -- \code{basic} & 8 & 118 \\
    \code{SabreSwap} -- \code{lookahead} & 8 & 116 \\
    \code{SabreSwap} -- \code{decay} & 8 & 116 \\
    \code{mini\_router} (tự viết, heuristic \code{basic}) & 10 & 113 \\
    \bottomrule
  \end{tabular}
\end{table}

Bảng~\ref{tab:routers} cho hai quan sát. Thứ nhất, \code{BasicSwap} tốn gấp 1.75 lần SABRE: nhìn
trước rất quan trọng. Thứ hai, \code{mini\_router} cài \emph{đúng} công thức~\eqref{eq:hbasic} vẫn
tốn 10 SWAP so với 8 của \code{SabreSwap basic}, vì bản cài đặt thật còn lọc ứng viên kỹ hơn và chạy
nhiều thử nghiệm. Hệ quả phương pháp: \textbf{so sánh một agent với baseline tự cài là vô nghĩa} --
baseline phải là \code{SabreSwap} của Qiskit. \code{LookaheadSwap} cho ít SWAP nhất ở đây nhưng chậm
khoảng 200 lần, nên không dùng được ở quy mô lớn.

\paragraph{Không phải SWAP nào cũng tốn chi phí.} Pass \code{ElidePermutations} xoá các SWAP viết
tường minh trong mạch gốc bằng cách đổi nhãn qubit, và \code{OptimizeSwapBeforeMeasure} xoá SWAP đứng
ngay trước phép đo bằng cách đổi nhãn bit cổ điển. Khi đếm SWAP để làm reward, cần loại các SWAP
``miễn phí'' này, nếu không agent bị phạt oan.

%--------------------------------------------------------------
\subsection{Kiểm chứng kết quả biên dịch}
\label{sec:verify}

Một router có thể sai theo hai cách \emph{độc lập}, và cả hai đều \textbf{im lặng} -- không có
exception nào:
\begin{enumerate}[noitemsep]
  \item \textbf{Không hợp lệ}: còn cổng hai qubit nằm ngoài coupling map. Kiểm tra bằng pass
    \code{CheckMap} (ghi \code{property\_set["is\_swap\_mapped"]}).
  \item \textbf{Không đúng}: unitary khác mạch gốc. Kiểm tra bằng so sánh toán tử -- nhưng phải bù
    lại hoán vị qubit do layout và SWAP gây ra.
\end{enumerate}

\begin{warnbox}[Ba bẫy kiểm chứng gặp phải khi thực nghiệm]
\textbf{(a)} \code{Operator(compiled).equiv(Operator(goc))} trả về \code{False} \emph{dù mạch đúng},
vì \code{Operator(compiled)} đọc theo qubit vật lý còn mạch gốc theo qubit logic. Phải dùng
\code{Operator.from\_circuit(compiled)}, hàm này đọc \code{compiled.layout} để bù hoán vị.\\[2pt]
\textbf{(b)} Khi tự route và tự theo dõi $\pi$, phải bù bằng \textbf{nghịch đảo} $\pi^{-1}$
(\code{inv[vật\_lý] = logic}). Dùng chính $\pi$ chỉ cho kết quả đúng khi hoán vị là một phép đổi
chỗ đơn lẻ -- đúng là trường hợp mà test nhỏ hay rơi vào, nên lỗi này dễ lọt qua kiểm thử.\\[2pt]
\textbf{(c)} \code{SabreLayout} \emph{tự route bên trong}. Nối nó với \code{ApplyLayout} theo công
thức chuẩn sẽ áp layout hai lần và sinh mạch không hợp lệ. Dấu hiệu nhận biết: pass nào ghi
\code{property\_set["final\_layout"]} thì đã route rồi.
\end{warnbox}

%--------------------------------------------------------------
\subsection{Synthesis: từ unitary tới chuỗi cổng}
\label{sec:synthesis}

Layout và routing lo chuyện \emph{ở đâu}; synthesis lo chuyện \emph{cổng gì}: cho một unitary $U$,
tìm chuỗi cổng ngắn nhất trong basis gates thực hiện được $U$. Có hai bài toán khác nhau về bản chất:
\emph{exact synthesis} trên gate set liên tục (máy NISQ) và \emph{approximate synthesis} trên gate set
rời rạc (máy chịu lỗi).

\paragraph{Exact synthesis một qubit.} Mọi $U\in U(2)$ viết được dưới dạng phân tích Euler
\begin{equation}
  U = e^{i\alpha}\,R_Z(\phi)\,R_Y(\theta)\,R_Z(\lambda).
\end{equation}
Ba góc khớp với ba bậc tự do của $SU(2)$, nên ba phép quay là tối thiểu và luôn đủ. Trên phần cứng
siêu dẫn, $R_Z$ gần như miễn phí (chỉ đổi pha của xung điều khiển) còn \code{sx} thì tốn, nên basis
\code{ZSX} dùng 3 \code{rz} và 2 \code{sx}; mọi phân tích đều chính xác tới sai số làm tròn
$\sim10^{-16}$.

\paragraph{Exact synthesis hai qubit -- định lý KAK.} Mọi $U\in SU(4)$ viết được thành
\cite{khaneja2001cartan, vatan2004optimal}
\begin{equation}
  U = (A_1\otimes A_2)\; \exp\!\big(i(c_1\,XX + c_2\,YY + c_3\,ZZ)\big)\; (B_1\otimes B_2),
\end{equation}
với $A_i, B_i$ là unitary một qubit. Ba hệ số $(c_1,c_2,c_3)$ -- \textbf{toạ độ Weyl} -- quyết định
hoàn toàn số cổng CX cần thiết, và chỉ có bốn khả năng (Bảng~\ref{tab:kak}). Nghĩa là
\textbf{3 CX luôn đủ} cho mọi unitary hai qubit \cite{vatan2004optimal}, và SWAP -- cái giá của routing
-- thuộc lớp khó nhất.

\begin{table}[h]
  \centering
  \caption{Định lý KAK: số CX tối thiểu cho unitary hai qubit (đã kiểm chứng bằng
  \code{TwoQubitBasisDecomposer}, sai số $\sim10^{-15}$).}
  \label{tab:kak}
  \small
  \begin{tabular}{@{}cl@{}}
    \toprule
    \textbf{Số CX} & \textbf{Loại unitary} \\
    \midrule
    0 & tích tensor $A\otimes B$, không tạo entanglement \\
    1 & CX, CZ, controlled-phase \\
    2 & iSWAP, $\sqrt{\mathrm{SWAP}}$ \\
    3 & SWAP và \emph{hầu hết} unitary ngẫu nhiên (tập cần ít hơn có độ đo 0) \\
    \bottomrule
  \end{tabular}
\end{table}

KAK còn là một chiến lược tối ưu: gom một khối cổng hai qubit liền nhau thành \emph{một} unitary rồi
tổng hợp lại. Một khối có 7 CX mà toạ độ Weyl chỉ cần 2 thì tiết kiệm được 5 cổng -- điều mà các pass
huỷ cổng từng cặp không bao giờ tìm ra. Chuỗi pass \code{Collect2qBlocks} $\to$
\code{ConsolidateBlocks} $\to$ \code{UnitarySynthesis} làm đúng việc này, và trên mạch thử nghiệm đã
giảm tổng số cổng \textbf{36\%} (144 $\to$ 92 cổng).

\paragraph{Approximate synthesis -- Clifford+T và T-count.} Trên máy tính lượng tử chịu lỗi, mạch
chạy trên qubit logic của một mã sửa lỗi, và mã đó chỉ thực hiện được một tập cổng rời rạc, thường là
Clifford+T:
\begin{itemize}[noitemsep]
  \item Nhóm \textbf{Clifford} ($H$, $S$, CX) có cài đặt chịu lỗi rẻ. Nhưng theo định lý
    \textbf{Gottesman--Knill} \cite{gottesman1998}, mạch chỉ gồm Clifford mô phỏng được hiệu quả bằng
    máy tính cổ điển -- nên Clifford một mình không thể có lợi thế lượng tử.
  \item Thêm cổng $T$ thì tập trở thành universal. Nhưng $T$ phải thực hiện qua \textbf{magic state
    distillation} \cite{bravyi2005magic}, tốn rất nhiều qubit vật lý và chu kỳ cho \emph{mỗi} cổng.
\end{itemize}
Vì vậy trên máy chịu lỗi, hàm cost chính là \textbf{T-count}. Định lý \textbf{Solovay--Kitaev}
\cite{dawson2006sk} bảo đảm mọi unitary một qubit xấp xỉ được tới sai số $\epsilon$ bằng
$O(\log^c(1/\epsilon))$ cổng, nhưng hằng số rất tệ. Thuật toán của \textbf{Ross--Selinger}
\cite{ross2016gridsynth} (plugin \code{gridsynth} trong Qiskit) xấp xỉ phép quay $R_Z$ với T-count
khoảng $3\log_2(1/\epsilon) + O(\log\log(1/\epsilon))$, tối ưu khi không dùng ancilla. Bảng~\ref{tab:tcount}
so sánh hai phương pháp \emph{ở cùng mức sai số thực tế} -- điều kiện cần cho một so sánh công bằng.

\begin{table}[h]
  \centering
  \caption{T-count khi xấp xỉ cùng một phép quay một qubit bằng Clifford+T. Sai số là khoảng cách
  toán tử \emph{đo thực tế} trên mạch tổng hợp ra.}
  \label{tab:tcount}
  \small
  \begin{tabular}{@{}lccc@{}}
    \toprule
    \textbf{Phương pháp} & \textbf{T-count} & \textbf{Tổng số cổng} & \textbf{Sai số} \\
    \midrule
    \code{gridsynth}, $\epsilon=10^{-1}$ & 18 & 60 & $9.4\times10^{-2}$ \\
    Solovay--Kitaev, bậc đệ quy 1 & 30 & 56 & $4.0\times10^{-2}$ \\
    Solovay--Kitaev, bậc đệ quy 2 & 120 & 216 & $2.5\times10^{-2}$ \\
    \code{gridsynth}, $\epsilon=10^{-2}$ & 40 & 104 & $7.3\times10^{-3}$ \\
    Solovay--Kitaev, bậc đệ quy 3 & 478 & 874 & $2.5\times10^{-3}$ \\
    \code{gridsynth}, $\epsilon=10^{-3}$ & \textbf{66} & 170 & $\mathbf{9.7\times10^{-4}}$ \\
    \code{gridsynth}, $\epsilon=10^{-4}$ & 84 & 220 & $8.6\times10^{-5}$ \\
    \midrule
    Solovay--Kitaev bậc 3 + \code{OptimizeCliffordT} & 316 & 747 & $2.5\times10^{-3}$ \\
    \bottomrule
  \end{tabular}
\end{table}

Ở sai số khoảng $2.5\times10^{-3}$, Solovay--Kitaev cần 478 cổng $T$ trong khi \code{gridsynth} chỉ
cần 66 cổng $T$ mà còn đạt sai số tốt hơn -- \textbf{tiết kiệm hơn 7 lần}. Sau khi tổng hợp, mạch
Clifford+T vẫn còn dư: các cổng $T$ và Clifford giao hoán theo những luật nhất định, nên đẩy các cổng
qua nhau có thể làm hai cổng $T$ gặp nhau ($T\cdot T = S$, mà $S$ là Clifford). Pass
\code{OptimizeCliffordT} khai thác điều này và giảm T-count \textbf{34\%} mà giữ nguyên sai số. Tìm
chuỗi rewrite tốt nhất trong không gian các luật như vậy là một bài toán tìm kiếm -- chính là hướng
của AlphaTensor-Quantum \cite{ruiz2025alphatensor} và các phương pháp RL trên ZX-calculus
\cite{riu2025zx}.

\begin{warnbox}[Bẫy: tham số bị bỏ qua im lặng]
Plugin synthesis gọi qua \code{transpile(..., unitary\_synthesis\_method="gridsynth")} chỉ tác dụng
lên \code{UnitaryGate}. Nếu đầu vào là một cổng \code{rz} thường, \code{BasisTranslator} dịch nó bằng
luật cố định \emph{trước khi} plugin được gọi, và tham số \code{epsilon} bị bỏ qua: T-count giữ
nguyên 127 cho mọi $\epsilon$ từ $10^{-1}$ tới $10^{-6}$. Bảng số liệu trông hợp lý nhưng vô nghĩa.
Cách tự bảo vệ: \textbf{luôn đo sai số thực} của mạch tổng hợp ra thay vì tin tham số đầu vào.
\end{warnbox}

%--------------------------------------------------------------
\subsection{Tự viết pass và tích hợp vào pipeline}
\label{sec:custompass}

Một phương pháp biên dịch mới -- kể cả một RL agent -- muốn chạy trong Qiskit thì cuối cùng phải đóng
gói thành một \code{TransformationPass}. Đoạn mã~\ref{lst:pass} là một pass hoàn chỉnh huỷ hai SWAP
liên tiếp trên cùng cặp qubit (đúng vì $\mathrm{SWAP}^2 = I$). Hai chi tiết dễ sai: phải so sánh
\emph{tập} qubit (\code{swap(1,2)} và \code{swap(2,1)} là cùng một cổng), và hai SWAP phải kề nhau
\emph{trên DAG} trên cả hai qubit, không chỉ gần nhau trong danh sách lệnh.

\begin{lstlisting}[caption={Một \code{TransformationPass} huỷ các cặp SWAP kề nhau trên DAG.}, label={lst:pass}]
from qiskit.transpiler.basepasses import TransformationPass
from qiskit.dagcircuit import DAGOpNode

class CancelAdjacentSwaps(TransformationPass):
    def run(self, dag):
        for node in list(dag.topological_op_nodes()):
            if node.name != "swap" or node not in dag.op_nodes():
                continue
            succ = [s for s in dag.successors(node) if isinstance(s, DAGOpNode)]
            # the next gate on BOTH wires must be the same SWAP node
            if (len(succ) == 1 and succ[0].name == "swap"
                    and set(succ[0].qargs) == set(node.qargs)
                    and len(list(dag.successors(node))) == 1):
                dag.remove_op_node(succ[0])
                dag.remove_op_node(node)
        return dag
\end{lstlisting}

Có hai cách gắn pass vào pipeline: dựng một \code{PassManager([...])} riêng, hoặc \textbf{thay một
stage} của preset pass manager, ví dụ \code{pm.routing = PassManager([MyRouter(cm)])}. Cách thứ hai
giữ nguyên toàn bộ phần còn lại (translation, optimization, scheduling) và chỉ đổi phần cần nghiên cứu
-- đây là cách một RL router sẽ được tích hợp. Pass mới luôn phải được kiểm tra tính đúng bằng
\code{Operator}: \emph{một pass sai luôn ``thắng'' benchmark}.
